% Chapter 0
%\chapter*{}
%\label{Chapter0}

%----------------------------------------------------------------------------------------
\newcommand{\keyword}[1]{\textbf{#1}}
\newcommand{\tabhead}[1]{\textbf{#1}}
\newcommand{\code}[1]{\texttt{#1}}
\newcommand{\file}[1]{\texttt{\bfseries#1}}
\newcommand{\option}[1]{\texttt{\itshape#1}}

%----------------------------------------------------------------------------------------
\clearpage
%\chapter*{Abstract}
\chapter*{\vspace*{-2em} Abstract}
\addcontentsline{toc}{chapter}{Abstract}
\vspace{-1.2cm}
\textbf{The impact of computational structural descriptions of PFAS compounds on the performance of classification models estimating their half-lives}

\vspace{0.2cm}

Per- and polyfluoroalkyl substances (PFAS) are persistent, bioaccumulative chemicals of growing concern due to their resistance to degradation and widespread presence in the environment. This master’s thesis investigates using density functional theory (DFT)-derived molecular descriptors combined with machine learning to predict PFAS persistence in terms of biological half-lives. Our objectives are to evaluate how the choice of quantum chemical method (specifically functional and basis set) influences descriptor reliability, to identify cost-effective computational strategies, and to analyze the impact of individual descriptors on classification performance.

Using three DFT functionals (B3LYP, M06-2X, and ωB97X-D) and three basis sets (6-31G(d,p), 6-31++G(d,p), and cc-pVTZ), we studied nine representative PFAS compounds. Key electronic descriptors, such as the HOMO–LUMO gap, chemical hardness, chemical potential, and electrophilicity, were calculated and shown to be particularly sensitive to method choice.

Three machine learning algorithms—k-nearest neighbors (kNN), decision tree (DT), and extreme gradient boosting (XGBoost)—were trained using descriptor sets from various DFT methods, supplemented with cheminformatics features and two biological descriptors, species and sex. The obtained qualitative structure–property relationship ([Q]SPR) models classified compounds into short-, medium-, or long-half-life categories. The best-performing model, XGBoost with ωB97X-D/cc-pVTZ descriptors, achieved approximately 90\% accuracy, while mid-level methods yielded similar results with significantly lower computational cost.

These results show that DFT-derived descriptors improve PFAS persistence classification and that method selection significantly affects descriptor quality. We propose a hybrid workflow integrating cheminformatics-based screening with targeted DFT descriptor computation as a scalable, efficient approach for analyzing large PFAS libraries. This work contributes to developing in silico tools for regulatory risk assessment and environmental pollutant prioritization.

\vspace{0.1cm}
\textbf{Keywords:} PFAS, persistence, half-life, density functional theory, molecular descriptors, chemoinformatics, [Q]SAR/[Q]SPR, machine learning, classification algorithms

\clearpage
%\chapter*{Streszczenie}
\chapter*{\vspace*{-2em} Streszczenie}
\vspace{-1.2cm}
\addcontentsline{toc}{chapter}{Streszczenie}

\textbf{Wpływ dokładności opisu struktury związków PFAS metodami chemii komputerowej na jakość modeli klasyfikacyjnych stosowanych do szacowania ich okresu półtrwania}

\vspace{0.2cm}

Substancje per- i polifluoroalkilowe (PFAS) to trwałe, bioakumulujące się substancje chemiczne,
które budzą coraz większe obawy ze względu na ich odporność na degradację i powszechną obecność w
środowisku. Niniejsza praca magisterska bada wykorzystanie deskryptorów molekularnych pochodzących z teorii funkcjonału gęstości (DFT) w połączeniu z uczeniem maszynowym w celu przewidywania trwałości PFAS w kategoriach biologicznych okresów półtrwania. Naszymi celami jest ocena, w jaki sposób wybór metody chemii kwantowej
(w szczególności funkcjonału i bazy) wpływa na niezawodność deskryptorów, identyfikacja opłacalnych strategii obliczeniowych i analiza wpływu poszczególnych deskryptorów na
wydajność klasyfikacji.

Używając trzech funkcjonałów DFT (B3LYP, M06-2X i ωB97X-D) i trzech zestawów bazowych (6-
31G(d,p), 6-31++G(d,p) i cc-pVTZ), zbadaliśmy dziewięć reprezentatywnych związków PFAS.
Obliczono kluczowe deskryptory elektroniczne, takie jak luka HOMO–LUMO, twardość chemiczna, potencjał chemiczny i elektrofilowość, i wykazano, że są one szczególnie wrażliwe na wybór metody.

Trzy algorytmy uczenia maszynowego — k najbliższych sąsiadów (kNN), drzewo decyzyjne (DT) i ekstremalne wzmocnienie gradientu (XGBoost) — zostały przeszkolone przy użyciu zestawów deskryptorów z różnych
metod DFT, uzupełnionych o cechy chemiczno-informatyczne i dwa deskryptory biologiczne,
gatunek i płeć. Uzyskane jakościowe modele relacji struktura–właściwość ([Q]SPR)
klasyfikowały związki do kategorii krótkiego, średniego lub długiego okresu półtrwania. Najlepiej działający
model, XGBoost z deskryptorami ωB97X-D/cc-pVTZ, osiągnął około 90\% dokładności, podczas gdy metody średniego poziomu dały podobne wyniki przy znacznie niższym
kosztach obliczeniowych.

Wyniki te pokazują, że deskryptory pochodzące z DFT poprawiają klasyfikację trwałości PFAS
a wybór metody obliczeń znacząco wpływa na jakość używanych deskryptorów. Proponujemy hybrydowy przepływ pracy integrujący przesiewanie oparte na chemioinformatyce z ukierunkowanymi obliczeniami deskryptora DFT
jako skalowalne, wydajne podejście do analizy dużych bibliotek PFAS. Niniejsza praca przyczynia się do
rozwoju narzędzi in silico do oceny ryzyka regulacyjnego i ustalania priorytetów zanieczyszczeń środowiska.

\vspace{0.1cm}

\textbf{Słowa kluczowe:} PFAS, trwałość, okres półtrwania, teoria funkcjonału gęstości, deskryptory molekularne, chemoinformatyka, [Q]SAR/[Q]SPR, uczenie maszynowe, algorytmy klasyfikacyjne
